% Page budget: 1.55 pages | Owns: augmentation topology, additional states, encoder, and closed-loop objective.
% WRITING GUIDE
% Purpose: Define the final augmentation architecture and explain how it is trained inside the known feedback loop.
% Start here: Current implementation decisions D-129, D-140, D-141, D-180, D-220, D-233, and D-237 in docs/decisions.md.
% Supporting material: Quinten's 18 September outline feedback, the Meeting-audit synthesis, Thesis-writeup/Documentation/TRAINING-DESIGN.md, and docs/writeup figures.
% Mandatory papers: Read Hoekstra's 2025 LFR augmentation paper, the 2026 first-principles augmentation paper, the encoder-initialization paper, and Kessels' Chapter 5 before writing this section.
% Research-plan anchor: Preserve Aspect 2's dynamic parallel augmentation objective; document the departure from open-loop training and earlier state routing.
% Current decisions: Separate physical and additional states, keep the controller outside the model, score the complete training window without burn-in, and use one network for the physical correction and added-state update.
% Choices to justify: Final reported routing, added-state dimension, ANN width and depth, zero initialization, encoder history, closed-loop objective, rollout length, full-window scoring, and validation horizon. Use the configuration of the reported thesis run, not a superseded routing experiment.
% Horizon check: Treat encoder history, scored training window, joint-estimation horizon, and free-run validation horizon separately; relate candidates to relevant stable modes and test sensitivity to slower dynamics and free-integrator accumulation.
% Implementation audit: Read scripts/gantry/gantry_interconnect_dynamic.py, gantry_dynamic/{model,config,controller,data,training}.py, and model_augmentation/fit_systems/{interconnect,closed_loop,pre_encoder}.py.
% Reference check: Verify sources for parallel LFR augmentation, encoder initialization, closed-loop state extension, and any claimed architectural precedent against the original paper or thesis.
% Cautions: Earlier routing plans are superseded; distinguish the truth-only hidden absorber from learned states; closed-loop fit alone does not prove autonomous plant recovery.
% Figures: Absorber schematic, author's updated architecture, and thesis-owned common-reference loops. No residual panel or horizon figure.
% Section is finished when: Every signal has a defined frame and timing, the RK4 boundary is explicit, and the described architecture matches the final code.
%
% SOURCE MAP
% Hoekstra et al. 2025 EJC: dynamic-parallel topology, additional states, encoder-based truncated simulation.
% Hoekstra et al. 2026 Automatica submission: extended LFR formulation, baseline-preserving model/encoder initialization, joint identification context.
% Hoekstra et al. 2026 encoder paper: reconstructability-based W^b initialization; it does not identify W^a or test dynamic augmentation.
% Kessels 2025 thesis: direct closed-loop rollout, controller equations, and reconstruction of the machine controller state for each window.
% Drenth 2025 thesis: LPV-specific augmentation precedent; not the unrestricted MLP structure used here.
% This project: gantry specialization, residual-loop implementation, exact routing, RK4 boundary, and verification.
%
% MUST ESTABLISH (agreed with Dirk, 2026-10-06; step 2a of the thesis-section skill)
% Contribution delivered: the gantry realisation of S-DP augmentation (CT self-scheduled LPV-LFR baseline
%   inside a DT parallel model, joint estimation in the ten combinations) and its identification under
%   feedback. Not new: the S-DP topology (Hoekstra), augmentation of self-scheduled LPV-LFR with added
%   states (Drenth Eq. 5.2), closed-loop training with a known controller (Kessels, inspiration only).
% III-A Augmented model
%  1. Dynamic, not static: static augmentation adds no states (EJC Sec. 2, "flexible modes"); refitting
%     theta_base adds none either. Argue from the general case, no absorber.
%  1b. The augmented model is Hoekstra's LFR augmentation structure (EJC Eq. 4, 2026 Eq. 6) with a FIXED
%     interconnection (EJC Sec. 3.2) realising S-DP; its baseline block is Section II's self-scheduled
%     LFR with Delta(Y) = Y I6: an LFR inside an LFR, which is what the section title means.
%     Well-posedness: acyclic graph (EJC Thm. 1, 2026 Cond. 6) plus Section II's M(Y) condition.
%  2. Parallel, not series. LEAD REASON (Dirk, 2026-10-06): the orthogonal-by-construction method that
%     Section IV builds on is defined for the additive (parallel) structure (Gyorok 2026 abstract, Sec. 1;
%     D-003), so parallel lets us use that existing work. Supporting: Hoekstra gives no universal rule,
%     the structure follows prior knowledge (2026 Sec. 1, 6.3); the baseline stays the explicit nominal
%     term (research plan Aspect 2); parallel needs only a feedforward network for baseline-preserving
%     init, series needs ResNets (EJC Sec. 4.3, 2026 Sec. 5.4, 6.3).
%  3. State augmentation only, h_aug = 0: output P^T q is exact kinematics (cf. Kessels Remark 5.3). TODO reason.
%  4. Correction on all physical rows and added states: S-DP definition (EJC), X and Y reachable (D-103);
%     marginal-mode risk on free axes stated. TODO: why correct position rows rather than keep kinematics.
%     Context: Hoekstra 2026 Sec. 7 (F1Tenth) detaches the free integrators and learns only the
%     input-to-velocity map because the full state is measured; here only positions are measured.
%     One network for f_aug and g_aug is an implementation statement (D-180), no scientific claim.
%  5. DT correction after the CT RK4 step (input held, M(Y) per stage, affine normalisation). TODO reason
%     versus force injection in the CT derivative.
%  6. Unrestricted nonlinear augmentation (depends on Y through the state) versus LPV-structured
%     augmentation (Drenth): departure from the research plan's "in LPV-LFR form". TODO reason. Acyclic
%     routing adds no instantaneous loops; baseline well-posedness is Section II's.
%  7. Zero output layer reproduces the baseline from the same physical initial state (Hoekstra 2026
%     Eq. 29, D-237); the split is not unique (Sec. 5.2), handed to Section IV. Added states are memory
%     and model order, defined up to a state transformation, no physical coordinates.
% III-B Encoder
%  8. Encoder needed: only positions measured; truncated windows for cost and stability (Beintema).
%  9. W^b from the reconstructability map of a ZOH linearisation at an equilibrium (encoder paper
%     Eqs. 15-17, 30-32), not fitted (Hoekstra 2026 Eq. 30); W^a random. Claim only better starting
%     values and faster convergence, final accuracy comparable (encoder paper Sec. 4.4). The paper
%     assumes a stable baseline; ours is marginally stable (free X, Y), the map exists because the
%     positions make the linearisation observable. The paper treats co-estimating theta_base as out of
%     scope (footnote 1); we co-estimate. Precedent for linearising at zero: encoder paper Sec. 4.2
%     ("around the 0 point"), no reason given. TODOs: why Y = 0; disclose nominal parameters in W^b
%     versus detuned start; why model-based rather than data-based init (candidate: cost, Sec. 4.4);
%     Kaiming versus Xavier (todo only).
% 10. Encoder estimates the full state incl. measured positions (SUBNET). TODO: mention Kessels Remark 5.3?
% 11. Keep encoder history, scored window and validation horizon apart; values in Setup. The
%     joint-estimation sensitivity horizon belongs to Section IV.
% III-C Closed-loop identification
% 12. Known-controller, reference-driven output-error identification, inspired by Kessels (Eqs. 5.12,
%     5.13; Remark 5.6 open-loop attempt failed). Related to the indirect approach (Forssell and Ljung
%     Sec. 5.2) as a linear-theory motivation, not a guarantee; direct identification is not called
%     inherently biased. Free X and Y axes: the SUBNET consistency argument that Hoekstra inherits
%     (2026 Sec. 5.5) assumes an incrementally exponentially output stable data-generating system
%     (Beintema 2023 Condition 1); the open-loop gantry is not (free integrators), the loop with the
%     stabilising controller can be. Precondition, not a proof, for our nonlinear model. Same principle
%     as Kessels (closed-loop simulation of an unstable open-loop system), scoped to marginal stability.
% 13. Our formulation (ours, "we"): u_hat = u + C x^c + D(y - y_hat). Derivation by subtracting the
%     machine's controller equations, equal to the shared-r, shared-u_ff loop of the figure under a
%     compatible controller, signals and initial state (assumption; whether the data meet it is a Setup
%     or Results check: 4 kHz controller on recorded r - y versus recorded feedback force).
%     Main consequence: x^c_tau = 0 per window, no controller initialisation; consistent with the
%     encoder (model state and model controller both start on the machine's condition at tau). Contrast:
%     Kessels' direct form needs the machine controller state (Remark 5.4). Error before tau is not
%     carried in; the free-run validation does carry it.
% 14. Step order from no plant feedthrough: no algebraic loop, no added delay.
% 15. Every window sample scored (as Eq. 22); theta_base, theta_aug, eta joint, controller fixed.
%     Joint estimation framing (Dirk, 2026-10-06): the goal is to estimate theta_base jointly, in the ten
%     combinations. Obstacle from Jan's work: the split is not unique (2026 Sec. 5.2) and approximately
%     initialised parameters stay near their start (EJC Sec. 4.3, 2026 Sec. 6.3). Handover: Section IV
%     extends Jan's method with orthogonality so that the estimate of theta_base is UNIQUE (this
%     strength, not "recovers the true parameters"). State what interpretable means: the estimated
%     parameters keep their physical meaning; cite Gyorok 2026 Sec. 3 ("With non-unique theta_b, the
%     baseline model could lose its physical meaning"). Operational form only; the general definition
%     is the Introduction's (todo there). "Negation" and Kessels' definition (Sec. 5.3.1.3) belong to
%     Section IV; his Sec. 6.2 recommendation to the Introduction gap. Do NOT mention Jan's parameter regularisation anywhere.
%     Do NOT say explicitly that Section III's formulation is the unconstrained comparison arm; let
%     Section IV extend it naturally.
% Setup pattern for hyperparameters (2026 Sec. 7): honest basis per value (physical insight and
%     trial-and-error, line-search, inherited from earlier results).
% 16. Closed-loop fit can mask plant error; plant-domain checks and controller transfer in Results.
% Notation: u = P u_act as in Sections II and IV; no "(eom)" shorthand; controller matrices distinct from A_c(Y).
%
% DRAFT NOTES (cycle 16, 2026-10-09)
% PROPOSED MUST ESTABLISH for III-E (the agreed block above has no PS2 item; inferred from the README
%   ownership row III, source-map row III-D and the agreed PS2 sketch of the README Math standard):
%  17. The zero output layer gives the parameters acting through xbar no gradient at the start; plain training
%      then leaves the added states nearly unused (PS2-METHOD Sec. 2.2; THESIS-RESULTS steps 2, 3).
%  18. S1 runs alongside the minimisation of V and fits W^a (theta^a_enc) and a temporary head to the current
%      model's closed-loop residual; adapted from Hefny S1 (Sec. 2). S2 fits g_aug once on frozen encoder pairs.
%      S3 is the plain method, as Downey refines a 2SR initialisation by BPTT (Sec. 4.2).
%  19. Screen and stopping rules as the code states them (PS2Spec), without min and max update counts (sketch).
%  20. M_PS2 named; per phase what changes; claim no stronger than THESIS-RESULTS step 3.
% CORRECTIONS to the agreed block (code, log and sources win):
%   - item 9 "Kaiming versus Xavier (todo only)": the thesis run uses Xavier uniform, gain 1
%     (entry _THESIS_FIXED wa_encoder_init, D-239); stated, todo removed.
%   - item 12 "Forssell and Ljung Sec. 5.2": only the LiU report numbering is verified, so the citation
%     carries no section number (README row III-C).
%   - item 1b "2026 Cond. 6": not verified in this session; EJC Thm. 1 alone is cited.
%   - item 2 supporting reason "no universal rule (2026 Secs. 1, 6.3)": not in the citation log; left out.
%   - item 3 "TODO reason" for h_aug = 0: README row III-A gives Kessels Remark 5.3, stated in the prose.
%   - item 15 "Section IV makes the estimate UNIQUE": Section IV Proposition 1 gives uniqueness on its
%     tuples and to first order, so the handover says "under the conditions it states".
%   - Hoekstra 2026 Eq. (22) already estimates theta_base jointly (theta = col(theta_base, theta_aug,
%     theta_LFR, theta_enc), p. 8); the adaptation is the closed loop, sample tau in the encoder, and xi.
%   - README row III-D and citation-log row "hefny2015supervised Sec. 3": the S1A/S2 regressions are in
%     Sec. 2 (p. 4); Sec. 3 is Related Work. Cited as Sec. 2, log row appended.
% NOTATION (one meaning per symbol): Section IV's encoder parameters theta_enc and training coordinates xi
%   are kept, so the encoder history is zeta (old xi_tau) and eta is no longer the encoder (appendix eta_0);
%   the controller state is x^fb (old s clashed with the appendix scale s); tilde marks the physical state,
%   so the encoder correction is psi_nl (old tilde psi); M names the inertia, so the PS2 horizon is n + 1
%   (it equals the encoder output window); e is left to the servo error of Section V, so the S1 residual is
%   omega; L_b is the length, so the PS2 objectives are V_1, V_2; N counts tuples in Section IV, so PS2 counts
%   are n_c, n_w, n_s. xi_Delta keeps the Delta clash recorded in the Section II todo on m_Delta.
% LEFT OUT (reader test or sketch): optimiser schedules and step sizes of Adam, S1 and S2; S1/S2 minimum and
%   maximum update counts (sketch); the S2 best-weight restore and its defect (PS2-METHOD Sec. 4.3, mechanics);
%   the fallback mechanics beyond "S2 is skipped"; detached storage of targets; the head's zero output layer;
%   gradient checkpointing, compilation, float64; positive raw masses versus positive-definite inertia (D-229,
%   reports are in theta_base); Hoekstra's "no universal rule" (unverified); development history of the routing
%   (Caution and anchor: final routing only, todo in III-A); the window-length reason (D-220, owned by the
%   Section V hyperparameter table).
% TODOS RESOLVED from the restored file: VERIFY keys (records its own resolution); missing bib entries
%   gyorok2026obc and forssell1999revisited (added to refs.bib with CHECK notes); W^a Kaiming/Xavier (code:
%   Xavier); missing open-loop drift evidence (THESIS-RESULTS Sec. 5 plans it, pointer to Section VI);
%   h_aug = 0 reason (Kessels Remark 5.3). Kept and rewritten: zero-init attribution, position rows,
%   LPV structure, Y = 0, nominal parameters in W^b, model- versus data-based W^b, figure todos, encoder
%   bib version, controller-rate check, Kessels Remark 5.3 alternative.
% STALE LINES in other sections (not edited): 05_setup.tex line 56 "Free-run validation (referenced from
%   Section III-C)": now Section III-D (sec:aug_closed_loop); 05_setup.tex lines 57, 59 use varphi for the
%   combinations, Section III uses theta_base and xi; 04_obc.tex header line "03_augmentation.tex line 667"
%   refers to the previous draft; the handover now states uniqueness under Section IV's conditions.
% SYMBOLS FOR SECTION V (values in the hyperparameter table): T_s, n, n_f, stride of T, n_xbar, width n_h and depth
%   of phi_aug (n_h also the head width), width and depth of psi_nl, rho, theta_base,0 (correct or detuned), Adam batch and rate, polish window count,
%   validation interval, beta_c; PS2 (PS2Spec): |B_1|, n_c, n_w, beta_1, n_s, gamma_max, |S|, |B_2|, n_c2,
%   n_w2, beta_2.
% CODE CONFIRMED this session: entry _THESIS_FIXED and CFG (up_sample 1, fs_new 4000, wa_encoder_init
%   xavier_uniform_gain1, lbfgs, stride 10, na_nb 29); GD/model.py build_model (network on [x, u], all 6 + n_a
%   rows, physical rows added to the RK4 result, zero_init_feed_forward_nn, encoder from the Y = 0 ZOH
%   linearisation scaled by normalize_linear_ss_matrices); SYS/gantry_linearization.py (nominal gantry_ss
%   parameters, Bc P for actuator forces); FS/pre_encoder.py forward (offsets on the linear part, x_off on the
%   physical rows only); GD/data.py compute_normalization; FS/blocks.py _rk4, _deriv_with,
%   combinations_from_free, _m_diff_bound; FS/closed_loop.py rollout, ControllerBank folding, step order,
%   validation_error, closed_loop_free_run_rms_batch; FS/interconnect.py loss and the after_step seam;
%   FS/lbfgs_polish.py decide and meter_regression; GD/training.py _gantry_meters (combo_err against the true
%   parameters); FS/ps2_opening.py (PS2Spec, calibration_split, after_step, _s1_step, _residual, _fit_map).
%   Algebra checked: residual form by subtraction; folded controller B T_y^-1, T_u C, T_u D T_y^-1; centring
%   cancels in h^n_base because P^T [I 0] mu_x = mu_y; tanh map gives (L_b^2/4) m_Delta^2 < rho^2 m_Sigma J_eff.
% LENGTH: a \draftfalse build in the scratchpad puts Section III from the bottom of the right column of p. 5 to
%   about 40 percent of the left column of p. 9, about 3.5 pages with two figures, against the 1.55-page budget
%   (2.3x). A second deletion pass (literature positioning, mechanics, sentences after displays) removed only the
%   line-search detail; the remaining sentences carry agreed MUST ESTABLISH items. The budget predates the
%   ownership row of D-244, which adds PS2, the optimiser, selection and the two named models. The required
%   content is twelve displays (model, normalisation, network, zero layer, encoder, W^b, two loop forms,
%   coordinates, problem, S1, S2) and the two planned figures. Proposed budget: 3.3 pages.
\section{Dynamic LPV-LFR Augmentation}\label{sec:augmentation}
\ifdraft
\begin{itemize}
\item Goal from Section~II: a model that also captures omitted dynamics, identified under feedback; names III-A to III-E (style profile)
\item Contribution share: the second contribution of Section~I, in its wording; PS2 and the joint estimation widen it (?)
\end{itemize}
\fi

Given the baseline of Section~\ref{sec:baseline}, our goal is a model that also
captures dynamics the baseline omits, identified from records measured under
feedback.
To this end, we define the augmented model (Section~\ref{sec:aug_structure}),
the encoder that estimates the initial state of each training window
(Section~\ref{sec:aug_encoder}) and the closed-loop simulation of a window
(Section~\ref{sec:aug_loop}).
We then state the identification problem (Section~\ref{sec:aug_closed_loop})
and the initialisation of the added states (Section~\ref{sec:aug_ps2}).
This delivers the second contribution of Section~I: dynamic parallel LFR
augmentation extended to a closed-loop, self-scheduled state-space setting with
learned additional states.
\todo{Missing: contribution wording. The section also delivers the joint estimation of $\theta_{\mathrm{base}}$ in its ten combinations and the data-driven initialisation of the added states, which the Introduction does not name. Candidate: ``Extend dynamic parallel LFR augmentation to a closed-loop, self-scheduled state-space setting, with joint estimation of the identifiable physical combinations and additional states initialised from data'' (Sections~\ref{sec:aug_closed_loop}, \ref{sec:aug_ps2}).}

\subsection{Augmented model}\label{sec:aug_structure}
\ifdraft
\begin{itemize}
\item If the omitted effects are dynamic, the six baseline states cannot represent them; refitting $\theta_{\mathrm{base}}$ or a static augmentation adds no states (EJC Sec.~2)
\item Hoekstra's LFR augmentation structure with a fixed interconnection; its baseline block is the LFR of Section~II, an LFR inside an LFR (EJC Secs.~3.1, 3.2; adopted)
\item Dynamic parallel because the OBC of Section~IV is defined for the additive structure; the baseline stays an explicit term (Gy\"or\"ok 2026 abstract; 2026 Table~1; D-003)
\item Output map not augmented: it is the exact map to the measured positions (Kessels Remark~5.3)
\item $f_{\mathrm{base}}$ is one RK4 step, input held, $M(Y)$ per stage, so self-scheduled; input $\uact$, output $\qact$ (code; this work)
\item The only algebraic loop is the scheduling loop of the baseline, well posed by \eqref{eq:wellposed} (EJC Thm.~1; Section~II)
\item Affine normalisation with training-record statistics, unlike Hoekstra Sec.~5.3; reason (?) (code; D-119)
\item $f_{\mathrm{aug}}$ writes all six physical rows, $X$ and $Y$ without stiffness; position rows (?); no LPV structure, unlike Drenth Eq.~(5.2) (?) (EJC Sec.~2; D-103)
\item One tanh MLP supplies $f^{\mathrm n}_{\mathrm{aug}}$ and $g^{\mathrm n}_{\mathrm{aug}}$; feedforward suffices for a parallel structure; added states have no physical coordinates (EJC Sec.~4.3; D-180)
\item Zero output layer reproduces the baseline from the same physical initial state, $\bar x=0$ after one step (2026 Eq.~(29); D-237; attribution (?))
\end{itemize}
\fi

If the effects that the baseline omits are dynamic, its six states cannot
represent them.
Refitting $\theta_{\mathrm{base}}$ adds no states.
Neither does a static augmentation, which corrects the state update of the
baseline without extending its state~\cite[Sec.~2]{hoekstra2025lfr}.
The augmentation therefore needs states of its own.

We use the LFR-based augmentation structure of Hoekstra et al., in which an
interconnection matrix connects a baseline block and a learning
block~\cite[Sec.~3.1]{hoekstra2025lfr}.
Our baseline block is the self-scheduled LFR of Section~\ref{sec:lfr}, so the
augmented model is an LFR whose baseline block is itself an LFR in $Y$.
We fix the interconnection, as they do~\cite[Sec.~3.2]{hoekstra2025lfr}, to the
dynamic-parallel structure with added states~\cite[Table~1]{hoekstra2026lfr}.
We choose the parallel structure because the orthogonal-by-construction
parametrisation that Section~\ref{sec:obc} extends is defined for additive
augmentation~\cite{gyorok2026obc}.
It also keeps the baseline as an explicit term of the state update.
The output map is not augmented.
It is the exact kinematic map to the measured positions, for which Kessels
finds no need of an output augmentation~\cite[Remark~5.3]{kessels2025ai}.

The baseline block enters through its one-step map $f_{\mathrm{base}}$, one
fourth-order Runge-Kutta (RK4) step of \eqref{eq:eom} over the sample period
$T_s$ with the input held over the step.
Each stage evaluates $M(Y)$ at its own state, so $f_{\mathrm{base}}$ remains
self-scheduled.
In the recorded signals, the augmented model is
(Fig.~\ref{fig:aug_structure})
\begin{equation}
\begin{aligned}
 \tilde x_{k+1}
 &=f_{\mathrm{base}}(\theta_{\mathrm{base}},\tilde x_k,u_{\mathrm{act},k})
   +f_{\mathrm{aug}}(\theta_{\mathrm{aug}},\hat x_k,u_{\mathrm{act},k}),\\
 \bar x_{k+1}
 &=g_{\mathrm{aug}}(\theta_{\mathrm{aug}},\hat x_k,u_{\mathrm{act},k}),\\
 \hat q_{\mathrm{act},k}&=h_{\mathrm{base}}(\tilde x_k)=P^\top[I_3\;\,0]\,\tilde x_k,
\end{aligned}
\label{eq:aug_transition}
\end{equation}
where $\tilde x=\col(q,\dot q)\in\R^6$ is the physical state of
\eqref{eq:lfr_G}, $\bar x\in\R^{n_{\bar x}}$ the added state,
$\hat x=\col(\tilde x,\bar x)$, $\uact\in\R^3$ the actuator force, which
$f_{\mathrm{base}}$ maps to $u=P\uact$ by \eqref{eq:Ptransform},
$\hat q_{\mathrm{act}}\in\R^3$ the predicted actuator positions,
$\theta_{\mathrm{base}}\in\R^{10}$ the combinations \eqref{eq:phi},
$\theta_{\mathrm{aug}}\in\R^{n_{\mathrm{aug}}}$ the parameters of the learning
block, $f_{\mathrm{base}}:\R^{10}\times\R^6\times\R^3\to\R^6$,
$f_{\mathrm{aug}}:\R^{n_{\mathrm{aug}}}\times\R^{6+n_{\bar x}}\times\R^3\to\R^6$
and
$g_{\mathrm{aug}}:\R^{n_{\mathrm{aug}}}\times\R^{6+n_{\bar x}}\times\R^3\to\R^{n_{\bar x}}$.
No learned map feeds the baseline block within a step, so the only algebraic
loop~\cite[Thm.~1]{hoekstra2025lfr} is the scheduling loop of the baseline,
well posed by \eqref{eq:wellposed}.

\begin{figure}[t]
 \centering
 \resizebox{\columnwidth}{!}{\input{figures/tex/augmentation_structure}}
 \caption{Augmented model. The physics step $f_{\mathrm{base}}$ and the
 network $\phi_{\mathrm{aug}}$ act in parallel on $\hat x_k$;
 $\phi_{\mathrm{aug}}$ supplies $f_{\mathrm{aug}}$ and $g_{\mathrm{aug}}$,
 and the output map is not learned ($h_{\mathrm{aug}}=0$). The encoder
 $\psi$ sets $\hat x$ at each window start. Here $u_k$ is the actuator force,
 during training the closed-loop input.
 \todo{No setpoint generator currently for the inputs. Too cramped; a two-column figure is a candidate.}
 }
 \label{fig:aug_structure}
\end{figure}

The model is identified in normalised coordinates.
Unlike Hoekstra et al., who only scale and take the state scale from a
baseline simulation~\cite[Sec.~5.3]{hoekstra2026lfr}, we centre and scale every
signal with statistics of the training records:
\begin{equation}
\begin{aligned}
 \tilde x^{\mathrm n}&=T_x(\tilde x-\mu_x),\\
 u^{\mathrm n}&=T_u(\uact-\mu_u),\qquad
 y^{\mathrm n}=T_y(\qact-\mu_y),
\end{aligned}
 \label{eq:aug_norm}
\end{equation}
where $T_\bullet=\operatorname{diag}(\sigma_\bullet)^{-1}$, $\mu_\bullet$ and
$\sigma_\bullet$ are the mean and standard deviation of each channel over the
training records, the state statistics are those of $q=P^{-\top}\qact$ and its
first difference over $T_s$, and $\bar x$ keeps its coordinates.
In these coordinates, the maps of \eqref{eq:aug_transition} become
$f^{\mathrm n}_{\mathrm{base}}=T_x(f_{\mathrm{base}}-\mu_x)$,
$f^{\mathrm n}_{\mathrm{aug}}=T_xf_{\mathrm{aug}}$,
$g^{\mathrm n}_{\mathrm{aug}}=g_{\mathrm{aug}}$ and
$h^{\mathrm n}_{\mathrm{base}}(\tilde x^{\mathrm n})=T_yP^\top[I_3\;\,0]\,T_x^{-1}\tilde x^{\mathrm n}$,
each evaluated at $\tilde x=T_x^{-1}\tilde x^{\mathrm n}+\mu_x$ and
$\uact=T_u^{-1}u^{\mathrm n}+\mu_u$.
\todo{Missing: reason for centring and for state statistics from the data, unlike~\cite[Sec.~5.3]{hoekstra2026lfr}. Candidate: their normalisation assumes zero-mean signals~\cite[Sec.~5.3]{hoekstra2026lfr}, which the recorded positions are not, and the encoder and the baseline block must share one state frame, which the data statistics give (D-119) (own reading).}

The learning block writes every physical row and depends on $Y$ without LPV
structure.
As in the dynamic-parallel structure, $f_{\mathrm{aug}}$ corrects the entire
physical state~\cite[Sec.~2]{hoekstra2025lfr}, so the learned block acts on the
coupled response of all three axes.
On $X$ and $Y$, a correction meets no stiffness
(\eqref{eq:damping_stiffness_matrices}), so a persistent correction accumulates
in position.
The learning block depends on $Y$ through $\hat x$ but, unlike the augmented
LPV-LFR of Drenth~\cite[Eq.~(5.2)]{drenth2025thesis}, has no LPV structure.
\todo{Missing: why the position rows are corrected rather than kept as the integrals of the velocities. Context: Hoekstra et al.\ detach the free integrators of their vehicle and learn only the input-to-velocity map, which needs measured velocities~\cite[Secs.~7.2, 7.4]{hoekstra2026lfr}; here only positions are measured. D-103 requires only that the learned block reaches $X$ and $Y$.}
\todo{Missing: reason for an unrestricted network rather than an LPV-structured augmentation. The research plan (Aspect~2) proposed extending the structure to the LPV setting following Drenth, and D-017 asks for an LPV-structured augmentation; no later decision supersedes it, so this is a departure. Candidate: none in the decisions.}

One multilayer perceptron (MLP) with tanh hidden layers and a linear output
layer supplies both learned maps in the coordinates of \eqref{eq:aug_norm}.
A plain feedforward network suffices, because the baseline is augmented
additively~\cite[Sec.~4.3]{hoekstra2025lfr}:
\begin{equation}
 \begin{bmatrix}
  f^{\mathrm n}_{\mathrm{aug}}(\theta_{\mathrm{aug}},\hat x^{\mathrm n},u^{\mathrm n})\\
  g^{\mathrm n}_{\mathrm{aug}}(\theta_{\mathrm{aug}},\hat x^{\mathrm n},u^{\mathrm n})
 \end{bmatrix}
 =\phi_{\mathrm{aug}}\bigl(\theta_{\mathrm{aug}},\col(\hat x^{\mathrm n},u^{\mathrm n})\bigr),
 \label{eq:aug_network}
\end{equation}
where $\hat x^{\mathrm n}=\col(\tilde x^{\mathrm n},\bar x)$,
$\phi_{\mathrm{aug}}:\R^{n_{\mathrm{aug}}}\times\R^{9+n_{\bar x}}\to\R^{6+n_{\bar x}}$
is the MLP, and its first six outputs form $f^{\mathrm n}_{\mathrm{aug}}$ and
the last $n_{\bar x}$ form $g^{\mathrm n}_{\mathrm{aug}}$.
The added states thus have no assigned physical coordinates: an invertible
transformation of $\bar x$, absorbed by the learned maps, leaves the
input-output behaviour unchanged.

Training starts from the baseline, with the output layer of
$\phi_{\mathrm{aug}}$ at zero:
\begin{equation}
 W_{\mathrm o}=0
 \quad\Longrightarrow\quad
 f^{\mathrm n}_{\mathrm{aug}}(\theta_{\mathrm{aug}},\cdot)=0,\quad
 g^{\mathrm n}_{\mathrm{aug}}(\theta_{\mathrm{aug}},\cdot)=0,
 \label{eq:aug_zero_init}
\end{equation}
where $W_{\mathrm o}\in\R^{(6+n_{\bar x})\times(n_{\mathrm h}+1)}$, a block of
$\theta_{\mathrm{aug}}$, collects the weight and bias of the output layer and
$n_{\mathrm h}$ is the hidden-layer width.
From the same physical initial state, the augmented model then reproduces the
baseline~\cite[Eq.~(29)]{hoekstra2026lfr}, with $\bar x_k=0$ after the first
step.
\todo{Missing: attribution of the all-zero output layer. It follows the public dynamic-parallel example code of Hoekstra et al.\ (D-237), while~\cite[Sec.~5.4.3]{hoekstra2026lfr} also permits random entries in the matrices the baseline behaviour does not need. Candidate: ``as in the public implementation of~\cite{hoekstra2026lfr}'', once Jan confirms which initialisation produced the reported dynamic-parallel results.}

\subsection{Encoder and initial state}\label{sec:aug_encoder}
\ifdraft
\begin{itemize}
\item Each window starts from a state of which only the positions are measured; an encoder estimates it from the history; short windows bound the cost and stabilise the optimisation (Beintema Secs.~2, 3)
\item Linear map plus nonlinear correction; the history includes sample $\tau$; the linear part acts on the scaled, uncentred history (encoder paper Eq.~(8), Eqs.~(11) to (17); 2026 Eq.~(22e); code)
\item $W^a$ Xavier as in 2026 Eq.~(31); correction output layer zero (D-239; code)
\item All states are estimated, the measured positions included (?: Kessels Remark~5.3)
\item $W^b$ from the reconstructability map of the ZOH linearisation at rest at $Y=0$ with the nominal parameters, computed rather than fitted as in 2026 Eq.~(30) (encoder paper Sec.~3.1; code)
\item The map exists although the baseline is marginally stable: the positions make the linearisation observable (encoder paper Secs.~2, 3.1; this work)
\item Better starting values and faster convergence, comparable final accuracy (encoder paper Sec.~4.4)
\item Open: why $Y=0$; nominal parameters in $W^b$; model- versus data-based (?)
\end{itemize}
\fi

Each training window starts from a state of which only the positions are
measured.
Following Beintema et al.~\cite[Secs.~2, 3]{beintema2023subnet}, an encoder
estimates this state from the input-output history up to the window start.
Simulating short windows from estimated states keeps the cost of a gradient
step independent of the record length and stabilises the optimisation.

As in the encoder of Hoekstra et al., the estimate is a linear map of the
history plus a nonlinear correction~\cite[Eq.~(8)]{hoekstra2026encoder}.
Unlike their identification criterion, which reads the history up to sample
$\tau-1$~\cite[Eq.~(22e)]{hoekstra2026lfr}, the history includes sample $\tau$,
as their reconstructability map does~\cite[Eqs.~(11) to~(17)]{hoekstra2026encoder}:
{\small
\begin{equation}
\begin{aligned}
 \hat x^{\mathrm n}_{\tau|\tau}
 &=\psi(\theta_{\mathrm{enc}},\zeta_\tau)
 =\begin{bmatrix}W^bT_\zeta\zeta_\tau-T_x\mu_x\\ W^aT_\zeta\zeta_\tau\end{bmatrix}
  +\psi_{\mathrm{nl}}\bigl(T_\zeta(\zeta_\tau-\mu_\zeta)\bigr),\\
 \zeta_\tau
 &=\col\bigl(q_{\mathrm{act},\tau-n},\dots,q_{\mathrm{act},\tau},\,
   u_{\mathrm{act},\tau-n},\dots,u_{\mathrm{act},\tau}\bigr),
\end{aligned}
\label{eq:aug_encoder}
\end{equation}
}
where $\tau$ is the window start, $n$ the encoder history,
$T_\zeta=\operatorname{blkdiag}(I_{n+1}\otimes T_y,\,I_{n+1}\otimes T_u)$ and
$\mu_\zeta=\col(\mathbf 1_{n+1}\otimes\mu_y,\,\mathbf 1_{n+1}\otimes\mu_u)$
apply \eqref{eq:aug_norm} to $\zeta_\tau\in\R^{6(n+1)}$,
$W^b\in\R^{6\times6(n+1)}$ and $W^a\in\R^{n_{\bar x}\times6(n+1)}$ give the
physical and added states, $W^a$ starts from a Xavier uniform
draw~\cite[Eq.~(31)]{hoekstra2026lfr},
$\psi_{\mathrm{nl}}:\R^{6(n+1)}\to\R^{6+n_{\bar x}}$ is a tanh MLP whose
output layer starts at zero, and $\theta_{\mathrm{enc}}\in\R^{n_{\mathrm{enc}}}$
collects $W^b$, $W^a$ and the parameters of $\psi_{\mathrm{nl}}$, so that
$\psi:\R^{n_{\mathrm{enc}}}\times\R^{6(n+1)}\to\R^{6+n_{\bar x}}$.
The encoder estimates all states, the measured positions included.
\todo{Missing: whether to mention the alternative of initialising the positions from the measurement, which an accurate output map allows~\cite[Remark~5.3]{kessels2025ai}. Candidate: leave it out; the encoder paper estimates the full state.}
\todo{Align the figure's past-only encoder-history labels ($y^{k-1}_{k-n_a}$) with \eqref{eq:aug_encoder}, which uses sample $\tau$ and one history length $n$.}
\todo{Missing: published version of \texttt{hoekstra2026encoder} (arXiv:2602.13108v1, Hoekstra, Gy\"or\"ok, T\'oth, Schoukens). Candidate: keep the arXiv entry unless an IFAC version appears before submission.}

The physical rows $W^b$ start from the baseline, following the model-based
initialisation of Hoekstra et al.~\cite[Sec.~3.1]{hoekstra2026encoder}.
For zero input, every rest state with $\Theta=0$ is an equilibrium of
\eqref{eq:eom}, because $K$ in \eqref{eq:damping_stiffness_matrices} vanishes on
$X$ and $Y$.
We linearise \eqref{eq:eom} at rest at $Y=0$, with the nominal parameters of
Table~\ref{tab:nominal_parameters} and the actuator force as input.
Its zero-order-hold discretisation over $T_s$, scaled by $T_x$, $T_u$ and $T_y$
of \eqref{eq:aug_norm}, gives $(A_d,B_d,C_d)$.
This model is marginally stable, whereas the encoder paper assumes a stable
baseline~\cite[Sec.~2]{hoekstra2026encoder}.
\todo{Missing: why linearise at $Y=0$. The encoder paper also linearises ``around the 0 point''~\cite[Sec.~4.2]{hoekstra2026encoder} without a reason. Candidate: the middle of the symmetric stroke (own reading).}
\todo{Missing: disclosure. The nominal parameters of $W^b$ are also the true parameters of the simulated system, so the encoder starts from oracle information even when $\theta_{\mathrm{base}}$ starts detuned; the encoder paper leaves co-estimation of the baseline parameters out of scope~\cite[footnote~1]{hoekstra2026encoder}. Candidate: state this as a limitation in Section~\ref{sec:setup}, or build $W^b$ from $\theta_{\mathrm{base},0}$.}

Unlike~\cite[Eq.~(30)]{hoekstra2026lfr}, which fits the baseline encoder to
simulated baseline states, we compute the reconstructability map of
$(A_d,B_d,C_d)$ directly~\cite[Eqs.~(15) to~(17)]{hoekstra2026encoder}:
\begin{equation}
 W^b=\begin{bmatrix}
  A_d^nO_n^{+} & \;r_n-A_d^nO_n^{+}T_n
 \end{bmatrix},
 \label{eq:aug_wb}
\end{equation}
where $O_n\in\R^{3(n+1)\times6}$ is the observability matrix of $(A_d,C_d)$
over $n+1$ samples, of full column rank because the measured positions make
$(A_d,C_d)$ observable despite its marginal stability, $O_n^{+}$ its
pseudoinverse, $T_n\in\R^{3(n+1)\times3(n+1)}$ the Toeplitz matrix of the
Markov parameters, $r_n\in\R^{6\times3(n+1)}$ the input-to-state map over the
same samples, and the two blocks act on the position and force parts of
$T_\zeta\zeta_\tau$.
The model-based initialisation gives better starting values and faster
convergence than a random encoder, at comparable final
accuracy~\cite[Sec.~4.4]{hoekstra2026encoder}.
\todo{Missing: why the model-based rather than the data-based initialisation~\cite[Sec.~3.2]{hoekstra2026encoder}. Candidate: it needs only a pseudoinverse~\cite[Sec.~4.4]{hoekstra2026encoder}; the same section notes that the data-based one might be preferred for nonlinear baselines, so the choice needs Dirk's reason.}

\subsection{Closed-loop simulation}\label{sec:aug_loop}
\ifdraft
\begin{itemize}
\item Records measured under feedback, model must predict under feedback: windows simulated in closed loop, not driven by the recorded input as in 2026 Eq.~(22) (Kessels Eqs.~(5.12), (5.13), Remark~5.6)
\item Linear case: the indirect approach keeps an output-error model consistent; a motivation, not a guarantee (Forssell and Ljung)
\item $X$ and $Y$ free: the open-loop gantry violates the stability condition behind SUBNET consistency; the loop can meet it (Beintema Cond.~1; 2026 Sec.~5.5; Kessels Ch.~5)
\item Both loops share the reference, feedforward and controller, written in normalised coordinates (code; D-140)
\item The model output depends on the state only, which fixes the step order: no algebraic loop through $D_{\mathrm{fb}}$ (code; Kessels Remark~5.1)
\item Subtracting gives our residual form; $x^{\mathrm c}_\tau=0$ sets the model's controller to the machine's state; no reconstruction, unlike Kessels Remark~5.4 (this work; code)
\item Equal for a compatible controller, signals and initial state; controller discretised at $T_s$ (?) (D-141)
\item A closed-loop fit can mask plant error; changed-controller demonstration in Section~VI (THESIS-RESULTS step~6)
\end{itemize}
\fi

The records are measured under feedback, and the model must predict the
machine under feedback.
We therefore simulate each window in closed loop with the known controller,
rather than driving the model with the recorded input as
in~\cite[Eq.~(22)]{hoekstra2026lfr}.
Kessels trains augmented models in this
way~\cite[Eqs.~(5.12), (5.13)]{kessels2025ai}.
He also reports that identifying the open-loop model from his closed-loop data
failed~\cite[Remark~5.6]{kessels2025ai}.
For linear systems, fitting the closed-loop response to the reference with a
known controller is the indirect approach~\cite{forssell1999revisited}.
The reference is uncorrelated with the noise, so an output-error model stays
consistent.
For our nonlinear model, this is a motivation, not a guarantee.

The free axes of the baseline give a second reason.
$X$ and $Y$ have no stiffness in \eqref{eq:damping_stiffness_matrices}, so the
open-loop gantry is marginally stable.
A difference in initial velocity leaves a permanent position offset.
The consistency argument of the identification
method~\cite[Sec.~5.5]{hoekstra2026lfr} rests on that of the subspace encoder
method, which assumes
an incrementally exponentially output stable data-generating
system~\cite[Cond.~1]{beintema2023subnet}.
The open-loop gantry does not meet this condition, whereas the loop with the
stabilising controller can.
Kessels uses the same principle to update models of systems that are unstable
in open loop~\cite[Ch.~5]{kessels2025ai}.
For our nonlinear model, this makes the condition attainable but does not prove
consistency.
Section~\ref{sec:results} compares an open-loop and a closed-loop rollout on
one record.

The controller of the records is known, linear and time invariant.
The machine loop and the model loop are driven by the same reference and feedforward through this
controller and differ only in the output they feed back
(Fig.~\ref{fig:aug_closed_loop}).
In the coordinates of \eqref{eq:aug_norm}, with controller state
$x^{\mathrm{fb}}$ for the machine and $\hat x^{\mathrm{fb}}$ for the model,
they are
\begin{equation}
\begin{aligned}
 x^{\mathrm{fb}}_{k+1}&=A_{\mathrm{fb}}x^{\mathrm{fb}}_k+B_{\mathrm{fb}}(r^{\mathrm n}_k-y^{\mathrm n}_k),\\
 u^{\mathrm n}_k&=C_{\mathrm{fb}}x^{\mathrm{fb}}_k+D_{\mathrm{fb}}(r^{\mathrm n}_k-y^{\mathrm n}_k)+u^{\mathrm{ff,n}}_k,\\
 \hat x^{\mathrm{fb}}_{k+1}&=A_{\mathrm{fb}}\hat x^{\mathrm{fb}}_k+B_{\mathrm{fb}}(r^{\mathrm n}_k-\hat y^{\mathrm n}_k),\\
 \hat u^{\mathrm n}_k&=C_{\mathrm{fb}}\hat x^{\mathrm{fb}}_k+D_{\mathrm{fb}}(r^{\mathrm n}_k-\hat y^{\mathrm n}_k)+u^{\mathrm{ff,n}}_k,
\end{aligned}
\label{eq:aug_direct_loop}
\end{equation}
where $x^{\mathrm{fb}},\hat x^{\mathrm{fb}}\in\R^{n_{\mathrm{fb}}}$,
$r^{\mathrm n}=T_y(r-\mu_y)$ is the normalised reference in actuator
coordinates, $u^{\mathrm{ff,n}}$ the normalised feedforward, which carries the
centring $-T_u\mu_u$, $(A_{\mathrm{fb}},B_{\mathrm{fb}},C_{\mathrm{fb}},D_{\mathrm{fb}})$
the controller from position error to actuator force with $B_{\mathrm{fb}}$,
$C_{\mathrm{fb}}$, $D_{\mathrm{fb}}$ scaled to these coordinates by $T_y^{-1}$
and $T_u$, and $\hat y^{\mathrm n}$, $\hat u^{\mathrm n}$ the model's output and
input.

The model output depends on the state only, so the feedthrough $D_{\mathrm{fb}}$
closes no algebraic loop, as Kessels requires of his output
augmentation~\cite[Remark~5.1]{kessels2025ai}.
Each sample is therefore evaluated in a fixed order: $\hat y^{\mathrm n}_k$ from
$\tilde x^{\mathrm n}_k$, then $\hat u^{\mathrm n}_k$, then the model and
controller states advance.
We implement the model loop in a residual form, obtained by subtracting the
machine's equations from the model's, which removes $r^{\mathrm n}$ and
$u^{\mathrm{ff,n}}$:
\begin{equation}
\begin{aligned}
 x^{\mathrm c}_{k+1}&=A_{\mathrm{fb}}x^{\mathrm c}_k+B_{\mathrm{fb}}(y^{\mathrm n}_k-\hat y^{\mathrm n}_k),
 \qquad x^{\mathrm c}_\tau=0,\\
 \hat u^{\mathrm n}_k&=u^{\mathrm n}_k+C_{\mathrm{fb}}x^{\mathrm c}_k+D_{\mathrm{fb}}(y^{\mathrm n}_k-\hat y^{\mathrm n}_k),
\end{aligned}
\label{eq:aug_residual_loop}
\end{equation}
where $x^{\mathrm c}=\hat x^{\mathrm{fb}}-x^{\mathrm{fb}}\in\R^{n_{\mathrm{fb}}}$,
and $u^{\mathrm n}$, $y^{\mathrm n}$ are the recorded input and output.
Because $x^{\mathrm c}$ is a difference of controller states,
$x^{\mathrm c}_\tau=0$ sets the model's controller to the machine's state at
each window start.

\begin{figure}[t]
 \centering
 \includegraphics[width=\columnwidth]{closed_loop_same_reference.pdf}
 \caption{Recorded plant loop (top) and augmented-model loop (bottom),
 driven by one shared reference $r_k$ and feedforward $u_k^{\mathrm{ff}}$,
 with the same feedback controller. The model output is evaluated
 before the input advances its state to the next sample.}
 \label{fig:aug_closed_loop}
\end{figure}

Unlike Kessels' direct form, the residual form needs no reconstruction of the
machine's controller state from the reference, the output and the
controller~\cite[Remark~5.4]{kessels2025ai}.
Model error made before $\tau$ is not carried into the window.
The validation free run of Section~\ref{sec:aug_closed_loop} carries it.

The residual form equals the model loop of \eqref{eq:aug_direct_loop} when the
recorded input was produced by this controller from the same output samples
and initial state.
The model loop runs the controller discretised at $T_s$.
The records were generated at a higher controller rate, so the equality is
approximate.
\todo{Missing: a check of this assumption on the thesis records. Candidate: run the controller discretised at $T_s$ on the recorded $r-y$ and compare its feedback force with the recorded one in Section~\ref{sec:setup}; a rollout with $\hat y^{\mathrm n}=y^{\mathrm n}$ gives $\hat u^{\mathrm n}=u^{\mathrm n}$ by construction and checks nothing.}

A close fit of the closed-loop response does not by itself show that the
plant model is correct, because the controller acts on the output error and
can mask plant error.
Section~\ref{sec:results} therefore also demonstrates the model under a changed
controller.

\subsection{Identification problem}\label{sec:aug_closed_loop}
\ifdraft
\begin{itemize}
\item Need from Section~II-C: training coordinates that keep \eqref{eq:lfr_admissibility} at every iterate; nine logarithmic, $m_\Delta$ bounded (D-191, D-229; code; this work)
\item Every $\xi$ satisfies \eqref{eq:lfr_admissibility}, so $M(Y)\succ0$ without clamp or penalty (D-229; this work)
\item $\vartheta=(\xi,\theta_{\mathrm{aug}},\theta_{\mathrm{enc}})$ minimises the closed-loop output error over every window sample; adapted from 2026 Eq.~(22); Kessels keeps the FP parameters fixed (Eq.~(5.12))
\item Adam, then an unbounded L-BFGS polish, adapted from Drenth (bounded L-BFGS-B) (D-171; code)
\item Selection: closed-loop free run over whole validation records, length-weighted RMS in metres; the polish is kept only if it lowers it and the combination error against the true values does not rise (oracle (?)) (D-171; code)
\item $n$, $n_f$ and the validation horizon are separate quantities (header)
\item The split between baseline and network is not unique; interpretable means the estimate keeps its physical meaning, which needs a unique estimate; Section~IV (2026 Secs.~5.2, 6.3; Gy\"or\"ok Sec.~3)
\item $\mathcal M_{\mathrm U}$ named; input $\uact$, output $\qact$; per phase what changes; purpose (THESIS-RESULTS steps~2, 4; notation (?))
\end{itemize}
\fi

Identifying the model needs coordinates for $\theta_{\mathrm{base}}$ that keep
\eqref{eq:lfr_admissibility} at every iterate, as Section~\ref{sec:params}
requires.
We train the nine positive combinations in logarithmic coordinates relative to
the start $\theta_{\mathrm{base},0}$, which keeps them positive.
Positivity does not bound $m_\Delta$, so we map it into the admissible set:
\begin{equation}
\begin{aligned}
 \theta_{\mathrm{base},j}(\xi)&=\theta_{\mathrm{base},0,j}\,e^{\xi_j},\qquad j\neq8,\\
 m_\Delta(\xi)&=\rho\,\frac{2}{L_b}\sqrt{m_\Sigma J_{\mathrm{eff}}}\,\tanh(\eta_0+s\,\xi_\Delta),
\end{aligned}
\label{eq:mdelta_map}
\end{equation}
where $\xi\in\R^{10}$ are the training coordinates, $\xi_\Delta=\xi_8$ is the
entry of $m_\Delta$ in the order of \eqref{eq:phi},
$\theta_{\mathrm{base},0}\in\R^{10}$ is the start, $\rho\in(0,1)$ a margin, and
$\eta_0$, $s$ make $\xi=0$ reproduce the start with
$\partial m_\Delta/\partial\xi_\Delta=m_{\Delta,0}$
(Appendix~\ref{app:lfr-wellposed}).
Since $|\tanh|<1$ and $\rho<1$, every $\xi$ satisfies
\eqref{eq:lfr_admissibility}, so $M(Y)\succ0$ for every $Y$ at every iterate,
without a clamp or penalty.

The trained parameters are $\vartheta=(\xi,\theta_{\mathrm{aug}},\theta_{\mathrm{enc}})$.
Kessels' closed-loop criterion, in contrast, keeps the first-principles
parameters fixed~\cite[Eq.~(5.12)]{kessels2025ai}.
We adapt the truncated simulation criterion of Hoekstra et
al.~\cite[Eq.~(22)]{hoekstra2026lfr}, with each window simulated in the closed
loop of Section~\ref{sec:aug_loop}, sample $\tau$ in the encoder and
$\theta_{\mathrm{base}}$ in $\xi$:
{\small
\begin{equation}
\begin{aligned}
 &\min_{\vartheta}\;V(\vartheta)=\frac{1}{|\mathcal T|\,n_f\,n_y}
  \sum_{\tau\in\mathcal T}\sum_{\ell=0}^{n_f-1}
  \norm{y^{\mathrm n}_{\tau+\ell}-\hat y^{\mathrm n}_{\tau+\ell|\tau}}_2^2
  \quad\text{s.t.}\\
 &\tilde x^{\mathrm n}_{k+1|\tau}
  =f^{\mathrm n}_{\mathrm{base}}\bigl(\theta_{\mathrm{base}}(\xi),
   \tilde x^{\mathrm n}_{k|\tau},\hat u^{\mathrm n}_{k|\tau}\bigr)
  +f^{\mathrm n}_{\mathrm{aug}}\bigl(\theta_{\mathrm{aug}},
   \hat x^{\mathrm n}_{k|\tau},\hat u^{\mathrm n}_{k|\tau}\bigr),\\
 &\bar x_{k+1|\tau}
  =g^{\mathrm n}_{\mathrm{aug}}\bigl(\theta_{\mathrm{aug}},
   \hat x^{\mathrm n}_{k|\tau},\hat u^{\mathrm n}_{k|\tau}\bigr),\qquad
  \hat y^{\mathrm n}_{k|\tau}=h^{\mathrm n}_{\mathrm{base}}(\tilde x^{\mathrm n}_{k|\tau}),\\
 &\hat u^{\mathrm n}_{k|\tau}=u^{\mathrm n}_k+C_{\mathrm{fb}}x^{\mathrm c}_{k|\tau}
  +D_{\mathrm{fb}}\bigl(y^{\mathrm n}_k-\hat y^{\mathrm n}_{k|\tau}\bigr),\\
 &x^{\mathrm c}_{k+1|\tau}=A_{\mathrm{fb}}x^{\mathrm c}_{k|\tau}
  +B_{\mathrm{fb}}\bigl(y^{\mathrm n}_k-\hat y^{\mathrm n}_{k|\tau}\bigr),\qquad
  x^{\mathrm c}_{\tau|\tau}=0,\\
 &\hat x^{\mathrm n}_{\tau|\tau}=\psi(\theta_{\mathrm{enc}},\zeta_\tau),
\end{aligned}
\label{eq:aug_loss}
\end{equation}
}
where $\mathcal T$ is the set of window starts in the training records, $n_f$
the window length, $n_y=3$, $k=\tau,\dots,\tau+n_f-1$, and
$\theta_{\mathrm{base}}(\xi)$ is \eqref{eq:mdelta_map}.

We adapt the optimiser sequence of Drenth et al., Adam followed by the bounded
L-BFGS-B~\cite[Sec.~5]{drenth2025lpvlfr}.
Adam runs on minibatches of windows, and the L-BFGS polish follows on a fixed
batch of windows.
The polish is unbounded, because \eqref{eq:mdelta_map} keeps every iterate
admissible.

Every fixed number of Adam updates, the model is simulated in closed loop over
each whole validation record.
The encoder and $x^{\mathrm c}=0$ act only at the record start, so model error
accumulates over the record.
The selection quantity is
$V_{\mathrm{val}}=\sum_rN_r\,\mathrm{RMS}_r/\sum_rN_r$, where $\mathrm{RMS}_r$
is the RMS error in metres of the predicted actuator positions on validation
record $r$ after its encoder history, and $N_r$ the record length.
Training restores the Adam checkpoint with the lowest $V_{\mathrm{val}}$.
The polish result is kept only if it lowers $V_{\mathrm{val}}$ and raises the
RMS relative error of the ten combinations against their true values by at
most a fraction $\beta_{\mathrm c}$.
The encoder history $n$, the window length $n_f$ and the whole-record
validation horizon are thus separate quantities.
\todo{Missing: the polish acceptance compares the combinations with their true values (\texttt{combo\_err}, $m_\Delta$ scaled by the mean actuator mass), an oracle available in simulation only (D-171). Candidate: keep the rule for the simulation study and state it as an oracle in Section~\ref{sec:setup}, or accept the polish on $V_{\mathrm{val}}$ alone.}

The minimiser of \eqref{eq:aug_loss} does not determine the split between
baseline and network, because $f^{\mathrm n}_{\mathrm{aug}}$ can also learn
changes that another $\theta_{\mathrm{base}}$ would
produce~\cite[Sec.~5.2]{hoekstra2026lfr}.
In the experiments of Hoekstra et al., approximately initialised physical
parameters stay close to their initial values~\cite[Sec.~6.3]{hoekstra2026lfr}.
The network meanwhile learns dynamics that the baseline could represent.
With a non-unique $\theta_{\mathrm{base}}$, the baseline can lose its physical
meaning~\cite[Sec.~3]{gyorok2026obc}.
The estimated combinations are interpretable here if they keep that meaning,
which requires a unique estimate.
Section~\ref{sec:obc} extends the method with a construction that makes the
estimate unique under the conditions it states.

We call $\mathcal M_{\mathrm U}$ the model \eqref{eq:aug_transition},
\eqref{eq:aug_network} with the encoder \eqref{eq:aug_encoder}, identified by
\eqref{eq:aug_loss}.
Its input is the recorded actuator force $\uact$ and its output the actuator
positions $\qact$.
It starts from $\theta_{\mathrm{base},0}$, \eqref{eq:aug_zero_init} and the
encoder of Section~\ref{sec:aug_encoder}.
The controller and the normalisation stay fixed.
Adam changes every entry of $\vartheta$.
Selection then restores the best Adam checkpoint.
The polish changes every entry of $\vartheta$ again, under the acceptance rule
above.
Section~\ref{sec:results} uses $\mathcal M_{\mathrm U}$ with and without added
states to show what plain training of the augmentation captures.
\todo{Missing: model notation. Candidate: the calligraphic $\mathcal M_{\mathrm U}$ and $\mathcal M_{\mathrm{PS2}}$, since plain $M$ clashes with $M(Y)$ (README open conflict on model names).}

\subsection{Initialisation of the added states}\label{sec:aug_ps2}
\ifdraft
\begin{itemize}
\item The zero output layer gives the parameters acting through $\bar x$ no gradient at the start; plain training then leaves the added states nearly unused (PS2-METHOD Sec.~2.2; THESIS-RESULTS steps~2, 3)
\item S1, after every update of $V$: the encoder's added-state outputs and a temporary head predict the current model's closed-loop residual; adapted from Hefny's first stage (Hefny Sec.~2; code)
\item S1 trains only $\theta_{\mathrm p}$ and $\theta^a_{\mathrm{enc}}$; it ends on a plateau of the calibration unexplained fraction; a screen skips S2 (code \texttt{PS2Spec}; README sketch)
\item S2 once: one-step regression of $g^{\mathrm n}_{\mathrm{aug}}$ on frozen encoder pairs; it moves all of $\theta_{\mathrm{aug}}$, so $f^{\mathrm n}_{\mathrm{aug}}$ changes too; it ends on a calibration plateau (code)
\item The stages give the added states predictive information and an update that propagates it, without a sufficiency guarantee (PS2-METHOD Secs.~1, 6.2.1)
\item S3 continues the minimisation of $V$, as Downey refines a two-stage regression initialisation by BPTT (Downey Sec.~4.2)
\item $\mathcal M_{\mathrm{PS2}}$ named; per phase what changes; selection may precede S2 (?) (THESIS-RESULTS step~3; PS2-METHOD Sec.~7.6)
\end{itemize}
\fi

The zero output layer \eqref{eq:aug_zero_init} leaves every parameter that acts
only through $\bar x$ without gradient at the start, because $\bar x$ reaches
the output only through $f^{\mathrm n}_{\mathrm{aug}}$.
Training from this start leaves the added states nearly unused and most of the
error of the omitted dynamics in place (Section~\ref{sec:results}).

After every update of $V$, our first stage S1 fits the encoder's added-state
outputs and a temporary head $h_{\mathrm p}$ to the current model's closed-loop
residual over the next $n+1$ samples.
This adapts the first-stage regression of Hefny et al., whose target is a
feature of future observations~\cite[Sec.~2]{hefny2015supervised}:
{\small
\begin{equation}
\begin{aligned}
 &\min_{\theta_{\mathrm p},\,\theta^a_{\mathrm{enc}}}\;
  V_1=\frac{1}{|\mathcal B_1|(n+1)n_y}\sum_{\tau\in\mathcal B_1}
  \norm{h_{\mathrm p}(\theta_{\mathrm p},\bar x_{\tau|\tau})
   -\omega_\tau/\omega_{\mathrm{rms}}}_2^2,\\
 &\omega_\tau=\col\bigl(y^{\mathrm n}_{\tau+\ell}
   -\hat y^{\mathrm n}_{\tau+\ell|\tau}\bigr)_{\ell=0}^{n},\\
 &\omega_{\mathrm{rms}}^2=\frac{1}{|\mathcal B_1|(n+1)n_y}
  \sum_{\tau\in\mathcal B_1}\norm{\omega_\tau}_2^2,
\end{aligned}
\label{eq:ps2_s1}
\end{equation}
}
where $\mathcal B_1$ is a minibatch of window starts in the training records
outside a calibration set $\mathcal R_{\mathrm{cal}}$, which holds one training
record of every record class of Section~\ref{sec:setup} with at least three
records, $\bar x_{\tau|\tau}$ is the added-state part of \eqref{eq:aug_encoder},
$\hat y^{\mathrm n}$ follows from the constraints of \eqref{eq:aug_loss} at the
current $\vartheta$ and is held fixed during the update,
$h_{\mathrm p}:\R^{n_{\mathrm p}}\times\R^{n_{\bar x}}\to\R^{(n+1)n_y}$ is an
MLP with one tanh hidden layer of width $n_{\mathrm h}$ and
parameters $\theta_{\mathrm p}\in\R^{n_{\mathrm p}}$, and
$\theta^a_{\mathrm{enc}}$ collects $W^a$ and the added-state rows of the output
layer of $\psi_{\mathrm{nl}}$.
S1 changes only $\theta_{\mathrm p}$ and $\theta^a_{\mathrm{enc}}$.

Every $n_{\mathrm c}$ updates, S1 evaluates the unexplained fraction
$\gamma=\sum\norm{h_{\mathrm p}(\theta_{\mathrm p},\bar x_{\tau|\tau})-\omega_\tau/\omega_{\mathrm{rms}}}_2^2/\sum\norm{\omega_\tau/\omega_{\mathrm{rms}}}_2^2$
on fixed windows of $\mathcal R_{\mathrm{cal}}$.
S1 ends when the lowest $\gamma$ of the last $n_{\mathrm w}$ evaluations is less
than a fraction $\beta_1$ below the lowest earlier $\gamma$.
If $\gamma>\gamma_{\max}$ at update $n_{\mathrm s}$ or at the end of S1, S2 is
skipped.
\todo{Missing: the S1 stopping rule also has a cap of 350 updates (\texttt{PS2Spec.s1\_cap}), which, not the plateau, ended S1 in every logged run (PS2-METHOD Secs.~5, 10.4); the README sketch excludes update counts. Candidate: add ``or after a fixed number of updates'' to the stopping sentence, with its value in Section~\ref{sec:setup}.}

At the end of S1, stage S2 fits the deployed added-state update once, by a
one-step regression on encoder pairs:
{\small
\begin{equation}
\begin{aligned}
 &\min_{\theta_{\mathrm{aug}}}\;
 V_2=\frac{1}{|\mathcal B_2|\,n_{\bar x}\,v}\sum_{k\in\mathcal B_2}
  \norm{g^{\mathrm n}_{\mathrm{aug}}(\theta_{\mathrm{aug}},\hat x^{\mathrm n}_{k|k},u^{\mathrm n}_k)
  -\bar x_{k+1|k+1}}_2^2,\\
 &v=\frac{1}{n_{\bar x}}\sum_{j=1}^{n_{\bar x}}
  \operatorname{var}_{k\in\mathcal B_2}\bigl(\bar x_{j,k+1|k+1}\bigr),
\end{aligned}
 \label{eq:ps2_s2}
\end{equation}
}
where $\mathcal B_2$ is a minibatch of a fixed set $\mathcal S$ of samples drawn
from the training records outside $\mathcal R_{\mathrm{cal}}$,
$\hat x^{\mathrm n}_{k|k}$ and $\bar x_{k+1|k+1}$ are outputs of
\eqref{eq:aug_encoder} for the histories ending at $k$ and $k+1$, computed once
at the end of S1 and held fixed, $\bar x_j$ is the $j$th added state, and
$\operatorname{var}$ is the sample variance over $\mathcal B_2$.
S2 moves every entry of $\theta_{\mathrm{aug}}$, the shared hidden layers
included, so $f^{\mathrm n}_{\mathrm{aug}}$ changes with
$g^{\mathrm n}_{\mathrm{aug}}$.

Every $n_{\mathrm c,2}$ steps, $V_2$ is evaluated on a fixed set of samples from
$\mathcal R_{\mathrm{cal}}$.
S2 stops when the lowest of the last $n_{\mathrm w,2}$ values is less than a
fraction $\beta_2$ below the lowest earlier value.

Stage S3 discards the head and continues the minimisation of $V$ from these
parameters, as Downey et al.\ refine a two-stage regression initialisation by
backpropagation through time~\cite[Sec.~4.2]{downey2017psrnn}.
S3 thus starts from added states that carry information predicting the
current model's closed-loop residual and from an update that propagates it,
without a guarantee that this information suffices.

We call $\mathcal M_{\mathrm{PS2}}$ the model of $\mathcal M_{\mathrm U}$
identified with these three stages.
Its input and output are those of $\mathcal M_{\mathrm U}$.
During S1, the updates of $V$ change every entry of $\vartheta$.
The S1 updates change $\theta_{\mathrm p}$ and $\theta^a_{\mathrm{enc}}$.
S2 changes only $\theta_{\mathrm{aug}}$.
S3 and the polish change every entry of $\vartheta$, as for
$\mathcal M_{\mathrm U}$.
Selection restores the best validated Adam checkpoint, so it can also return
one saved before S2.
Section~\ref{sec:results} compares $\mathcal M_{\mathrm{PS2}}$ with
$\mathcal M_{\mathrm U}$, which differ only in how the added states start.
\todo{Missing: whether the selected checkpoint of each reported run lies after S2 (PS2-METHOD Sec.~7.6). Candidate: report the selected update number beside the S2 boundary in Section~\ref{sec:results}.}
